% !PRINTED 13
\weberpage{13}

%% [running head, not body text -- commented out 2026-08-07] Page furniture (running
%% head, folio, signature) is excluded from this edition by standing policy.
%% \gesperrt{\S.\,4.} \hfill \gesperrt{Gruppe der Nebengruppen.} \hfill 13

\gesperrt{2. Durch die beiden Gleichungen}

\weq{2}{AA = A,}
\weq{3}{AB = BA,}

\gesperrt{worin $A$ ein bestimmter, $B$ jeder beliebige Theil von $P$ sein kann, wird ausgesagt, dass $A$ ein Normaltheiler von $P$ ist.}

Denn nach (2) ist $A$ ein Theiler von $P$ und nach (3) ist für jedes Element $b$ von $P$
\[
b^{-1}Ab = A,
\]
d.\,h. $A$ ist ein Normaltheiler.

\gesperrt{3. Sind $A$ und $B$ zwei Theiler von $P$ (Gruppen), so ist eine der beiden Gleichungen}

\weq{4}{AB = A, \quad BA = A,}

\gesperrt{die nothwendige und hinreichende Bedingung dafür, dass $B$ ein Theiler von $A$ ist.}

Denn wenn $B$ eine Gruppe und ein Theiler der Gruppe $A$ ist, und wenn $a$ in $A$, $b$ in $B$ und also auch in $A$ enthalten ist, so ist auch $ab$ in $A$ enthalten. $A$ enthält also jedes Element von $AB$.

Weil aber zweitens $B$ als Gruppe auch das Element 1 enthält, so enthält $AB$ auch jedes Element von $A$, und also ist $AB$ mit $A$ identisch.

Ebenso sieht man, dass $BA$ mit $A$ identisch ist. Andererseits folgt aus jeder der Gleichungen (4), da $A$ als Gruppe das Einheitselement enthält, dass jedes Element von $B$ in $A$ enthalten, also $B$ ein Theiler von $A$ ist.

\gesperrt{4. Bei der Composition der Theile bildet das System der Nebengruppen zu einem Normaltheiler von $P$ selbst eine Gruppe, in der der Normaltheiler $N$ die Einheit bildet, und}
\[
Nb, \quad Nb^{-1}
\]
\gesperrt{entgegengesetzte Elemente sind.}

Denn ist $N$ ein Normaltheiler von $P$ und
\weq{5}{N, \; N_1 = Nb_1, \; N_2 = Nb_2, \dots}
das System der Nebengruppen, so ist
\[
N_h N_k = N b_h N b_k.
\]
